\section{11. Deep Dive: Scaling Laws for RL}
% ============================================================

\begin{bluf}
If the conversation goes deep on scaling, you need to reason about: how performance scales with compute, how problem difficulty changes the picture, and the three-axis tradeoff. This section builds on the Jones paper and connects to the RL scaling framework.
\end{bluf}

\subsection{The Core Scaling Relationships}

\textbf{For pretraining} (Kaplan/Chinchilla): Loss $\propto C^{-\alpha}$ where $C$ is compute. Power law. Straight line on log-log. Chinchilla says scale params and data equally.

\textbf{For RL} (Jones): Elo $\propto \log(C)$ in the active regime, with the curve shifting linearly with problem difficulty. This means:
\begin{itemize}
  \item Doubling compute gives a \textit{constant} Elo improvement (not proportional)
  \item Harder problems need \textit{exponentially} more compute for the same performance
  \item But the \textit{shape} of the curve is the same --- it just shifts right
\end{itemize}

\textbf{Key difference:} Pretraining scaling is power-law in loss (diminishing but never-ending). RL scaling is sigmoid-shaped --- there's a ``takeoff'' compute below which nothing happens, an active regime of rapid improvement, and a ceiling (perfect play). This matters because it means RL has \textit{phase transitions}: below a threshold, your compute is wasted.

\subsection{Train-Time vs.\ Test-Time Tradeoff}

From Jones: 10$\times$ more training $\approx$ 15$\times$ less inference search needed. This is log-linear.

\textbf{For LLMs:} This maps to: more RL training should reduce the chain-of-thought length / number of retries needed at inference. A well-trained model ``knows'' the answer; a poorly-trained model needs to ``think'' through it.

\textbf{Anthropic's framing:} The RL scaling team optimizes the Pareto frontier across three axes (pretraining, RL, test-time). The train/test tradeoff determines where the efficient frontier sits. More RL training = cheaper inference = better Claude Code economics.

\subsection{The Compute Allocation Problem}

Given a fixed total compute budget, how do you split between:
\begin{itemize}
  \item Bigger base model (pretraining)?
  \item More RL training on top?
  \item More test-time compute (longer chains of thought)?
\end{itemize}

This is an empirical question with no known closed-form answer. The framing: ``How much am I willing to burn compute vs.\ how much am I willing to burn dollars on people's time to give scaffolding or rewards?''

\begin{deeper}
The Jones paper shows this tradeoff is measurable in controlled settings. An open question for LLMs: does the same log-linear substitution hold? Or does the reward hacking problem mean that RL compute has diminishing returns that pretraining compute doesn't? This is an actual research question you could investigate.
\end{deeper}

% ============================================================
\section{12. Deep Dive: Hardware \& Pipeline}
% ============================================================

\begin{bluf}
Frontier labs run on heterogeneous hardware: TPUs, NVIDIA GPUs (Hopper/Blackwell), Trainium, and potentially AMD. The ``training pipeline held together by duct tape'' means massive wins from systems engineering. You need to reason about bottlenecks across hardware families and compare them through the intelligence-per-Watt lens.
\end{bluf}

\subsection{Why Hardware Understanding Matters for RL}

RL training is \textit{more} hardware-intensive than pretraining because:
\begin{itemize}
  \item You must \textbf{sample from the policy} during training (generate full completions) --- this is expensive inference interleaved with training
  \item You need \textbf{multiple models in memory}: policy + reference + reward model (+ critic for PPO). GRPO helps by dropping the critic, but you still need 3 models.
  \item The \textbf{batch size is determined by rollout count}: you need enough samples per prompt to estimate advantages reliably (GRPO uses groups of $G$ samples)
  \item \textbf{Communication overhead}: distributing rollout generation across many accelerators, then aggregating rewards and advantages
\end{itemize}

\subsection{The Key Bottlenecks}

\textbf{Memory:} Holding 3--4 large models simultaneously. This is why GRPO's critic elimination matters --- it's a 25--50\% memory reduction.

\textbf{Throughput:} Generation (autoregressive decoding) is memory-bandwidth-bound, not compute-bound. The KV cache grows linearly with sequence length. For long-context RL (30-hour coding sessions), this becomes the binding constraint.

\textbf{Communication:} In distributed training, gradient synchronization across nodes. For RL specifically: collecting rollouts from many parallel environments and broadcasting reward signals.

\textbf{Utilization:} GPUs/TPUs sit idle during parts of the RL loop (e.g., waiting for rollout completion before computing advantages). Pipeline parallelism and async rollout generation are active areas of optimization.

\subsection{The Three Hardware Families}

\begin{tabular}{@{}p{1.6cm}p{2.2cm}p{3.5cm}@{}}
\toprule
\textbf{Platform} & \textbf{RL Strengths} & \textbf{Key Differences} \\
\midrule
\textbf{TPU} (v5e/v7) & Systolic array, HBM3E, Pallas kernels & Pipeline-oriented. Memory: HBM $\to$ VMEM $\to$ VREGs. 1M+ Ironwood chips committed. \\
\textbf{NVIDIA} (H100/B200) & Tensor cores, CUDA ecosystem, Flash Attention & Thread-oriented. Blackwell: TMEM, FP4/FP6, 5$\times$ inference claim. \\
\textbf{AMD} (MI300X) & 192GB HBM3 (2.4$\times$ H100), 5.3 TB/s bandwidth & 64-wide wavefronts (vs 32-wide warps), LDS 64KB, MFMA units. HIP $\approx$ CUDA but ecosystem immature. \\
\bottomrule
\end{tabular}

\medskip

\textbf{Intelligence-per-Watt comparison:} For RL workloads, the key metric is \textit{useful FLOP/s per Watt on RL-relevant operations}, not peak. Since generation is memory-bandwidth-bound, platforms with higher HBM bandwidth (AMD MI300X: 5.3 TB/s vs H100: 3.35 TB/s) may deliver better RL intelligence per Watt despite lower peak FLOP/s. The generation bottleneck means Amdahl's law dominates: a 5$\times$ generation speedup (Blackwell's claim) has outsized impact on RL wall-clock time.

\subsection{Simon Boehm's Pedagogical Method}

Simon Boehm's pedagogical method:
\begin{enumerate}
  \item \textbf{Start naive} --- write the simplest correct implementation
  \item \textbf{Profile} --- use hardware profilers to identify the actual bottleneck (not the assumed one)
  \item \textbf{Fix one thing} --- apply exactly one optimization
  \item \textbf{Measure} --- quantify the improvement
  \item \textbf{Repeat} --- next bottleneck
  \item \textbf{Document honestly} --- including what didn't work and what you don't understand
\end{enumerate}

This is the cycle-time methodology applied to systems engineering. Each step is small, measurable, and produces a shareable result.

\begin{sayit}[Systems understanding]
``I think about RL training as a systems problem as much as an algorithmic one. The reason GRPO matters isn't the math --- it's that dropping the critic frees 50\% of GPU memory, which you can reallocate to larger rollout batches or a bigger policy model. Similarly, the reason long-context RL is hard isn't attention complexity --- it's KV cache memory growing linearly with sequence length while you're already holding 3 models. The wins are in understanding where the bottleneck \textit{actually} is, not where you assume it is.''
\end{sayit}

\subsection{MoE Memory Budget Analysis}

Dario confirmed Anthropic is running MoE models. Under MoE, expert weights compete directly with the KV cache for HBM bandwidth and capacity. For a hypothetical 70B MoE with 8 experts and 2 active per token: all 8 expert weight tensors must reside in HBM (only routing is sparse --- weights stay loaded). Total expert weight memory $\approx 8 \times \text{FFN params} \times 2$ bytes (BF16). On H100 (80 GB HBM): expert weights consume the majority of available memory, leaving a narrow budget for KV cache and activations. On MI300X (192 GB HBM): comfortable --- expert weights are resident with substantial headroom.

The remaining HBM after model weights determines maximum context length for RLVR rollouts:
\[
\text{max context} \propto \frac{\text{HBM} - \text{model weights}}{\text{KV cache per token}}
\]

\textbf{Intelligence-per-Watt framing:} The hardware that can hold more context generates longer reasoning chains per training step. Longer rollouts = richer verification signal = more intelligence per Watt from the RL budget. MI300X's 192 GB HBM is not just a comfort margin --- it directly expands the maximum rollout length and reasoning depth accessible to the RL optimizer.

\subsection{Dario's ``Big Blob'' and PE}

Dario described 7 load-bearing factors for model quality. Factors 6--7 (numerical stability and conditioning) are performance engineering problems: ``getting the big blob of compute to flow in a laminar way'' is a precision management requirement, not a throughput optimization. Factor 5 (``objective function that can scale to the moon'') maps directly to the verification thesis --- RLVR's deterministic rewards are what make the objective scalable.

\begin{sayit}[Dario's factors]
``The CEO described numerical stability as one of seven load-bearing factors for model quality. Precision management --- FP8/BF16/FP32 accumulation --- isn't just a throughput optimization, it's a correctness requirement. I've worked on this across three hardware platforms.''
\end{sayit}

% ============================================================
